\clearpage
\oldpage[II]{217}
\phantomsection
\addcontentsline{toc}{section}{Errata and Addenda (List 1)}
\section*{Errata and Addenda}
\begin{center}
\large (List 1)
\end{center}

\subsection*{Chapter 0}

\paragraph{(0, 2.1.3).}
On line 15 from the bottom of p.~21, replace ``inter-section'' by
``intersection''.

\paragraph{(0, 2.1.6).}
On line 13 of p.~22, replace
\[
  Z_i\cap C\left(\bigcup_{j\ne i}Z_j\right)
  \quad\hbox{by}\quad
  \bigcup_{j\ne i}Z_j.
\]

\paragraph{(0, 3.1.6).}
On line 17 from the bottom of p.~25, replace
$\Gamma(U,\sh{F})$ by $\Gamma(V,\sh{F})$.

\paragraph{(0, 4.1.1).}
On line 12 of p.~36, replace $\sh{G}\to\sh{F}$ by
$\sh{B}\to\sh{A}$.

\paragraph{(0, 4.2.1).}
On line 11 from the bottom of p.~39, replace ``sheaves'' by ``sheaf''.

\paragraph{(0, 5.1.2).}
After line 14 of p.~45, add:

If $\sh{G}$ is an $\sh{O}_Y$-module generated by its sections over $Y$,
then $f^*(\sh{G})$ is generated by its sections over $X$, since $f^*$ is a
right-exact functor.

\paragraph{(0, 5.2.4).}
On line 18 from the bottom of p.~46, replace $f_U$ by $f_U^*$.

\paragraph{(0, 5.3.4).}
Before line 16 from the bottom of p.~47, add:

If
\[
  \sh{A}\to\sh{B}\to\sh{C}\to\sh{D}\to\sh{E}
\]
is an exact sequence of $\sh{O}_X$-modules, and if
$\sh{A},\sh{B},\sh{D},\sh{E}$ are coherent, then $\sh{C}$ is coherent.

\paragraph{(0, 5.4.1).}
After line 2 of p.~49, add:

Suppose that $\sh{O}_X$ is coherent, and let $\sh{F}$ be a coherent
$\sh{O}_X$-module.  If, at a point $x\in X$, $\sh{F}_x$ is a free
$\sh{O}_x$-module of rank $n$, then there exists a neighbourhood $U$ of
$x$ such that $\sh{F}|U$ is locally free of rank $n$; indeed,
$\sh{F}_x$ is then isomorphic to $\sh{O}_x^n$, and the assertion follows
from (0, 5.2.7).

\paragraph{(0, 6.6.2).}
Replace lines 10 and 11 of p.~59 by:

Indeed, this follows from (0, 6.4.1, d).

\paragraph{(0, 6.7.1).}
Before line 10 from the bottom of p.~59, add:

If $f$ is a flat morphism, we also say that $X$ is \emph{flat over $Y$},
or \emph{$Y$-flat}.

\paragraph{(0, 7.2.9).}
On line 11 from the bottom of p.~65, replace $M_n$ by $M_{n-1}$.

\subsection*{Chapter I}

\paragraph{(I, 1.2.1).}
On line 18 of p.~83, replace $K$ by $k$ (twice).

\paragraph{(I, 1.2.7).}
On line 16 from the bottom of p.~84, replace $A$ by $A'$ (twice) and
$\rad$ by $\rad'$; on line 17 from the bottom, replace
$X$ by $X'$.

\paragraph{(I, 1.7.3) and (I, 2.2.4).}
The statements given under these numbers can, in accordance with a
remark of J.~Tate, be generalized as follows.

\subsection*{1.8. Morphisms from locally ringed spaces to affine schemes}
\label{ega2:addendum:subsection:I.1.8}

\begin{proposition}[1.8.1]
\label{ega2:addendum:I.1.8.1}
Let $(S,\sh{O}_S)$ be an affine scheme and $(X,\sh{O}_X)$ a locally
ringed space.  There is a canonical bijection from the set of ring
homomorphisms
\[
  \Gamma(S,\sh{O}_S)\to\Gamma(X,\sh{O}_X)
\]
to the set of morphisms of ringed spaces
$(\psi,\theta):(X,\sh{O}_X)\to(S,\sh{O}_S)$ such that, for every
$x\in X$, $\theta_x^\sharp:\sh{O}_{\psi(x)}\to\sh{O}_x$ is a local
homomorphism.
\end{proposition}

\oldpage[II]{218}
\begin{proof}
First note that, if $(X,\sh{O}_X)$ and $(S,\sh{O}_S)$ are any two
ringed spaces, a morphism $(\psi,\theta)$ from $(X,\sh{O}_X)$ to
$(S,\sh{O}_S)$ canonically defines a ring homomorphism
\[
  \Gamma(\theta):\Gamma(S,\sh{O}_S)\to\Gamma(X,\sh{O}_X),
\]
and hence a first map
\begin{equation}
\label{ega2:addendum:I.1.8.1.1}
  \rho:\Hom((X,\sh{O}_X),(S,\sh{O}_S))
  \to\Hom(\Gamma(S,\sh{O}_S),\Gamma(X,\sh{O}_X)).
\tag{1.8.1.1}
\end{equation}
Conversely, under the hypotheses of the proposition, put
$A=\Gamma(S,\sh{O}_S)$ and consider a ring homomorphism
$\vphi:A\to\Gamma(X,\sh{O}_X)$.  For every $x\in X$, the set of
$f\in A$ such that $\vphi(f)(x)=0$ is a prime ideal of $A$, since
$\sh{O}_x/\mathfrak{m}_x=\kres(x)$ is a field.  It is therefore an
element of $S=\Spec(A)$, which we again denote by ${}^a\vphi(x)$.
Moreover, for every $f\in A$, by definition (0, 5.5.2),
\[
  {}^a\vphi^{-1}(D(f))=X_f,
\]
which proves that ${}^a\vphi$ is a continuous map $X\to S$.

We next define a homomorphism of $\sh{O}_S$-modules
\[
  \widetilde{\vphi}:\sh{O}_S\to{}^a\vphi_*(\sh{O}_X).
\]
For every $f\in A$, one has
$\Gamma(D(f),\sh{O}_S)=A_f$ (I, 1.3.6); to $s/f\in A_f$ associate
\[
  (\vphi(s)|X_f)(\vphi(f)|X_f)^{-1}
  \in\Gamma(X_f,\sh{O}_X)
  =\Gamma(D(f),{}^a\vphi_*(\sh{O}_X)).
\]
Passing from $D(f)$ to $D(fg)$ shows at once that this defines a
homomorphism of sheaves and hence a morphism of ringed spaces
$({}^a\vphi,\widetilde{\vphi})$.  With the same notation, and writing
$y={}^a\vphi(x)$, one has (0, 3.7.1)
\[
  \widetilde{\vphi}_x^\sharp(s_y/f_y)
  =(\vphi(s)_x)(\vphi(f)_x)^{-1}.
\]
Since $s_y\in\mathfrak{m}_y$ is, by definition, equivalent to
$\vphi(s)_x\in\mathfrak{m}_x$, the homomorphism
$\widetilde{\vphi}_x^\sharp:\sh{O}_y\to\sh{O}_x$ is local.  We have
therefore defined a second map
\[
  \sigma:\Hom(\Gamma(S,\sh{O}_S),\Gamma(X,\sh{O}_X))
  \to\mathfrak{L},
\]
where $\mathfrak{L}$ denotes the set of morphisms
$(\psi,\theta):(X,\sh{O}_X)\to(S,\sh{O}_S)$ for which
$\theta_x^\sharp$ is local for every $x\in X$.

It remains to prove that $\sigma$ and the restriction of $\rho$ to
$\mathfrak{L}$ are inverses.  The definition of
$\widetilde{\vphi}$ immediately gives
$\Gamma(\widetilde{\vphi})=\vphi$, and hence
$\rho\circ\sigma$ is the identity.  To see that $\sigma\circ\rho$ is
the identity, begin with $(\psi,\theta)\in\mathfrak{L}$ and put
$\vphi=\Gamma(\theta)$.  The hypothesis on $\theta_x^\sharp$ gives,
after passage to residue fields, a monomorphism
\[
  \theta^x:\kres(\psi(x))\to\kres(x)
\]
such that, for every $f\in A=\Gamma(S,\sh{O}_S)$,
\[
  \theta^x(f(\psi(x)))=\vphi(f)(x).
\]
Thus $f(\psi(x))=0$ is equivalent to $\vphi(f)(x)=0$, which proves
that ${}^a\vphi=\psi$.  On the other hand, the definitions make the
diagram
\[
  \xymatrix{
    A\ar[r]^\vphi\ar[d] &
    \Gamma(X,\sh{O}_X)\ar[d]\\
    A_{\psi(x)}\ar[r]^{\theta_x^\sharp} &
    \sh{O}_x
  }
\]
commutative; the analogous diagram with $\theta_x^\sharp$ replaced
by $\widetilde{\vphi}_x^\sharp$ is also commutative.  It follows from
(0, 1.2.4) that
$\widetilde{\vphi}_x^\sharp=\theta_x^\sharp$, and consequently that
$\widetilde{\vphi}=\theta$.
\end{proof}

\begin{env}[1.8.2]
\label{ega2:addendum:I.1.8.2}
When $(X,\sh{O}_X)$ and $(Y,\sh{O}_Y)$ are locally ringed spaces, we
shall consider morphisms
$(\psi,\theta):(X,\sh{O}_X)\to(Y,\sh{O}_Y)$ such that, for every
$x\in X$, $\theta_x^\sharp:\sh{O}_{\psi(x)}\to\sh{O}_x$ is a local
homomorphism.

\oldpage[II]{219}
Whenever we speak henceforth of a \emph{morphism of locally ringed
spaces}, we shall always mean a morphism of this kind.  With this
definition of morphisms, locally ringed spaces form a category.  For
two objects $X,Y$ of this category, $\Hom(X,Y)$ will therefore denote
the set of morphisms of locally ringed spaces from $X$ to $Y$ (the set
denoted $\mathfrak{L}$ in \sref{I.1.8.1}).  When we wish to consider
the set of morphisms of ringed spaces from $X$ to $Y$, we shall write
$\Hom_{\mathrm{rs}}(X,Y)$ to avoid ambiguity.  The map
\sref{I.1.8.1.1} is thus written
\begin{equation}
\label{ega2:addendum:I.1.8.2.1}
  \rho:\Hom_{\mathrm{rs}}(X,Y)
  \to\Hom(\Gamma(Y,\sh{O}_Y),\Gamma(X,\sh{O}_X)),
\tag{1.8.2.1}
\end{equation}
and its restriction
\begin{equation}
\label{ega2:addendum:I.1.8.2.2}
  \rho':\Hom(X,Y)
  \to\Hom(\Gamma(Y,\sh{O}_Y),\Gamma(X,\sh{O}_X))
\tag{1.8.2.2}
\end{equation}
is functorial in $X$ and $Y$ in the category of locally ringed spaces.
\end{env}

\begin{corollary}[1.8.3]
\label{ega2:addendum:I.1.8.3}
Let $(Y,\sh{O}_Y)$ be a locally ringed space.  For $Y$ to be an affine
scheme, it is necessary and sufficient that, for every locally ringed
space $(X,\sh{O}_X)$, the map \sref{I.1.8.2.2} be bijective.
\end{corollary}

\begin{proof}
Proposition~\sref{I.1.8.1} shows that the condition is necessary.
Conversely, suppose it holds and put $A=\Gamma(Y,\sh{O}_Y)$.  The
hypothesis and \sref{I.1.8.1} show that the functors
\[
  X\longmapsto\Hom(X,Y)
  \quad\hbox{and}\quad
  X\longmapsto\Hom(X,\Spec(A))
\]
from locally ringed spaces to sets are isomorphic.  This implies the
existence of a canonical isomorphism $Y\to\Spec(A)$ (cf.~0, 8).
\end{proof}

\begin{env}[1.8.4]
\label{ega2:addendum:I.1.8.4}
Let $S=\Spec(A)$ be an affine scheme.  Denote by $(S',A')$ the ringed
space whose underlying space consists of one point and whose
structure sheaf $A'$ is the (necessarily simple) sheaf on $S'$ defined
by the ring $A$.  Let $\pi:S\to S'$ be the unique map.  For every open
subset $U$ of $S$, the canonical map
\[
  \Gamma(S',A')=\Gamma(S,\sh{O}_S)\to\Gamma(U,\sh{O}_S)
\]
defines a $\pi$-morphism $\iota:A'\to\sh{O}_S$ of sheaves of rings.
Thus there is a canonical morphism of ringed spaces
\[
  i=(\pi,\iota):(S,\sh{O}_S)\to(S',A').
\]
For every $A$-module $M$, denote by $M'$ the simple sheaf on $S'$
defined by $M$, which is evidently an $A'$-module.  It is clear that
$i_*(\widetilde M)=M'$ (I, 1.3.7).
\end{env}

\begin{lemma}[1.8.5]
\label{ega2:addendum:I.1.8.5}
With the notation of \sref{I.1.8.4}, for every $A$-module $M$, the
canonical functorial $\sh{O}_S$-homomorphism (0, 4.4.3.3)
\begin{equation}
\label{ega2:addendum:I.1.8.5.1}
  i^*(i_*(\widetilde M))\to\widetilde M
\tag{1.8.5.1}
\end{equation}
is an isomorphism.
\end{lemma}

\begin{proof}
Both sides of \sref{I.1.8.5.1} are right exact in $M$ (the functor
$M\mapsto i_*(\widetilde M)$ being evidently exact) and commute with
direct sums.  By viewing $M$ as the cokernel of a homomorphism
$A^{(I)}\to A^{(J)}$, one is reduced to the case $M=A$, where the
assertion is immediate.
\end{proof}

\begin{corollary}[1.8.6]
\label{ega2:addendum:I.1.8.6}
Let $(X,\sh{O}_X)$ be a ringed space and let $u:X\to S$ be a morphism
of ringed spaces.

\oldpage[II]{220}
For every $A$-module $M$, there is, with the notation of
\sref{I.1.8.4}, a canonical functorial isomorphism of
$\sh{O}_X$-modules
\begin{equation}
\label{ega2:addendum:I.1.8.6.1}
  u^*(\widetilde M)\isoto u^*(i^*(M')).
\tag{1.8.6.1}
\end{equation}
\end{corollary}

\begin{corollary}[1.8.7]
\label{ega2:addendum:I.1.8.7}
Under the hypotheses of \sref{I.1.8.6}, for every $A$-module $M$ and
every $\sh{O}_X$-module $\sh{F}$, there is a canonical isomorphism,
functorial in $M$ and $\sh{F}$,
\begin{equation}
\label{ega2:addendum:I.1.8.7.1}
  \Hom_{\sh{O}_S}(\widetilde M,u_*(\sh{F}))
  \isoto\Hom_A(M,\Gamma(X,\sh{F})).
\tag{1.8.7.1}
\end{equation}
\end{corollary}

\begin{proof}
By (0, 4.4.3) and Lemma~\sref{I.1.8.5}, there is a canonical
isomorphism of bifunctors
\[
  \Hom_{\sh{O}_S}(\widetilde M,u_*(\sh{F}))
  \isoto\Hom_{A'}(M',i_*(u_*(\sh{F}))).
\]
The right-hand side is plainly $\Hom_A(M,\Gamma(X,\sh{F}))$.  Under
the canonical homomorphism \sref{I.1.8.7.1}, an
$\sh{O}_S$-homomorphism
$h:\widetilde M\to u_*(\sh{F})$ (equivalently, a $u$-morphism
$\widetilde M\to\sh{F}$) corresponds to
\[
  \Gamma(h):M\to\Gamma(S,u_*(\sh{F}))=\Gamma(X,\sh{F}).
\]
\end{proof}

\begin{env}[1.8.8]
\label{ega2:addendum:I.1.8.8}
With the notation of \sref{I.1.8.4}, it is clear from (0, 4.1.1) that
every morphism of ringed spaces $(\psi,\theta):X\to S'$ is equivalent
to the data of a ring homomorphism
$A\to\Gamma(X,\sh{O}_X)$.  Proposition~\sref{I.1.8.1} may therefore
be interpreted as giving a canonical bijection
\[
  \Hom(X,S)\isoto\Hom(X,S'),
\]
where the right-hand side means morphisms of ringed spaces, since $A$
is not in general a local ring.

More generally, let $X,Y$ be locally ringed spaces and let $(Y',A')$
be the ringed space whose underlying space is one point and whose
sheaf of rings is the simple sheaf defined by
$\Gamma(Y,\sh{O}_Y)$.  Then \sref{I.1.8.2.1} may be interpreted as a
map
\begin{equation}
\label{ega2:addendum:I.1.8.8.1}
  \rho:\Hom_{\mathrm{rs}}(X,Y)\to\Hom(X,Y').
\tag{1.8.8.1}
\end{equation}
Corollary~\sref{I.1.8.3} says that affine schemes are characterized
among locally ringed spaces as those for which the restriction of
$\rho$ to $\Hom(X,Y)$,
\begin{equation}
\label{ega2:addendum:I.1.8.8.2}
  \rho':\Hom(X,Y)\to\Hom(X,Y'),
\tag{1.8.8.2}
\end{equation}
is bijective for every locally ringed space $X$.  In a later chapter
we shall generalize this definition, thereby associating to an
arbitrary ringed space $Z$ (and not only to one whose underlying
space is a point) a locally ringed space which we shall again call
$\Spec(Z)$.  This will be the starting point of a ``relative'' theory
of preschemes over arbitrary ringed spaces, extending the results of
Chapter~I.
\end{env}

\begin{env}[1.8.9]
\label{ega2:addendum:I.1.8.9}
The pairs $(X,\sh{F})$, where $X$ is a locally ringed space and
$\sh{F}$ an $\sh{O}_X$-module, form a category: a morphism
$(X,\sh{F})\to(Y,\sh{G})$ is a pair $(u,h)$ consisting of a morphism
of locally ringed spaces

\oldpage[II]{221}
$u:X\to Y$ and a $u$-morphism $h:\sh{G}\to\sh{F}$ of modules.  For
fixed $(X,\sh{F})$ and $(Y,\sh{G})$, these morphisms form a set denoted
$\Hom((X,\sh{F}),(Y,\sh{G}))$.  The assignment
$(u,h)\mapsto(\rho'(u),\Gamma(h))$ is a canonical map
\begin{equation}
\label{ega2:addendum:I.1.8.9.1}
\begin{split}
  \Hom((X,\sh{F}),(Y,\sh{G}))\to
  \Hom\bigl(
    &(\Gamma(Y,\sh{O}_Y),\Gamma(Y,\sh{G})),\\
    &(\Gamma(X,\sh{O}_X),\Gamma(X,\sh{F}))
  \bigr),
\end{split}
\tag{1.8.9.1}
\end{equation}
functorial in $(X,\sh{F})$ and $(Y,\sh{G})$; the right-hand side is
the set of di-homomorphisms corresponding to the rings and modules in
question (0, 1.0.2).
\end{env}

\begin{corollary}[1.8.10]
\label{ega2:addendum:I.1.8.10}
Let $Y$ be a locally ringed space and $\sh{G}$ an $\sh{O}_Y$-module.
For $Y$ to be an affine scheme and $\sh{G}$ a quasi-coherent
$\sh{O}_Y$-module, it is necessary and sufficient that, for every pair
$(X,\sh{F})$ consisting of a locally ringed space $X$ and an
$\sh{O}_X$-module $\sh{F}$, the canonical map \sref{I.1.8.9.1} be
bijective.
\end{corollary}

We leave to the reader the task of spelling out the proof, which is
modelled on that of \sref{I.1.8.3} and uses \sref{I.1.8.1} and
\sref{I.1.8.7}.

\begin{remark}[1.8.11]
\label{ega2:addendum:I.1.8.11}
The statements (I, 1.7.3), (I, 1.7.4), and (I, 2.2.4), as well as the
definition (I, 1.6.1), are special cases of \sref{I.1.8.1};
similarly, (I, 2.2.5) follows from \sref{I.1.8.7}.
Corollary~\sref{I.1.8.7} also gives (I, 1.6.3), and hence (I, 1.6.4),
as a special case.  Indeed, if $X$ is affine and
$\Gamma(X,\sh{F})=N$, the functors
\[
  M\longmapsto
  \Hom_{\sh{O}_S}(\widetilde M,u_*(\widetilde N))
  \quad\hbox{and}\quad
  M\longmapsto
  \Hom_{\sh{O}_S}(\widetilde M,(N_{[\vphi]})\supertilde)
\]
(where $\vphi:A\to\Gamma(X,\sh{O}_X)$ corresponds to $u$) are
isomorphic by \sref{I.1.8.7} and (I, 1.3.8).  Finally, (I, 1.6.5),
and hence (I, 1.6.6), follows from \sref{I.1.8.6} and from the fact
that, for every $f\in A$, the $A_f$-modules
\[
  N'\otimes_{A'}A_f
  \quad\hbox{and}\quad
  (N'\otimes_{A'}A)_f
\]
(with the notation of (I, 1.6.5)) are canonically isomorphic.
\end{remark}

\paragraph{Insert after (I, 3.2.8):}

\begin{env}[3.2.9]
\label{ega2:addendum:I.3.2.9}
It follows from \sref{I.1.8.1} that (I, 3.2.2) may be sharpened as
follows.  The prescheme
\[
  Z=\Spec(B\otimes_A C)
\]
is not only a product of $X=\Spec(B)$ and $Y=\Spec(C)$ in the category
of $S$-preschemes, but also in the category of locally ringed spaces
over $S$ (with the definition of $S$-morphisms modelled on
(I, 2.5.2)).  The proof of (I, 3.2.6) also shows that, for arbitrary
$S$-preschemes $X,Y$, the prescheme $X\times_S Y$ is not only their
product in the category of $S$-preschemes, but also in the category of
locally ringed spaces over the prescheme $S$.
\end{env}

\paragraph{(I, 4.4.5).}
On line 14 from the bottom of p.~126, replace $B$ by $A$; on line 13
from the bottom, replace ``$A$-algebra'' by ``$B$-algebra''.

\paragraph{(I, 5.3.5).}
On line 2 from the bottom of p.~132, replace $Y\to S$ by
$g:Y\to S$.

\paragraph{(I, 6.3.2.1).}
On line 6 from the bottom of p.~144, read
\[
  D(g_i)\subset W.
\]

\paragraph{(I, 7.3.8).}
Replace the present text, which is erroneous, by the following:

Let $X,Y$ be integral preschemes; then $\sh{R}(X)$ and $\sh{R}(Y)$
are respectively a quasi-coherent $\sh{O}_X$-module and a
quasi-coherent $\sh{O}_Y$-module (I, 7.3.3).  If
$f:X\to Y$ is a dominant morphism, there is a canonical homomorphism
of $\sh{O}_X$-modules
\begin{equation}
\label{ega2:addendum:I.7.3.8.1}
  \tau:f^*(\sh{R}(Y))\to\sh{R}(X).
\tag{7.3.8.1}
\end{equation}

\oldpage[II]{222}
First suppose that $X=\Spec(A)$ and $Y=\Spec(B)$ are affine, with
$A,B$ integral domains.  Then $f$ corresponds to an injective
homomorphism $B\to A$, which extends to a monomorphism
$L\to K$ from the fraction field $L$ of $B$ into the fraction field
$K$ of $A$.  The homomorphism \sref{I.7.3.8.1} then corresponds to the
canonical homomorphism
\[
  L\otimes_B A\to K
\]
of (I, 1.6.5).

In the general case, for every pair of nonempty affine open subsets
$U\subset X$, $V\subset Y$ such that $f(U)\subset V$, the preceding
construction defines a homomorphism $\tau_{U,V}$.  If
$U'\subset U$, $V'\subset V$, and $f(U')\subset V'$, then
$\tau_{U,V}$ extends $\tau_{U',V'}$; the assertion follows by gluing.
If $x,y$ are the generic points of $X,Y$, respectively, then
$f(x)=y$ and
\[
  \bigl(f^*(\sh{R}(Y))\bigr)_x
  =\sh{O}_y\otimes_{\sh{O}_y}\sh{O}_x
  =\sh{O}_x
\]
by (0, 4.3.1), so $\tau_x$ is an isomorphism.

\paragraph{(I, 9.5.2).}
In the last three lines of p.~176 and the first five lines of p.~177,
replace $X'$ everywhere by $Y'$ and $X''$ by $Y''$.

\paragraph{(I, 10.5.4).}
On line 11 from the bottom of p.~187, replace (10.3.4) by (10.3.5).

\paragraph{(I, 10.6.3).}
On line 3 from the bottom of p.~189, replace $X$ by $\mathfrak{X}$.

\paragraph{(I, 10.14.2).}
On line 5 of p.~210, replace ``coherent $\sh{O}_X$-module'' by
``coherent ideal of $\sh{O}_X$''; on line 7, replace ``coherent
$\sh{O}_X$-modules'' by ``coherent ideals of $\sh{O}_X$''.
